\subsubsection{Order Types and Instructions}

\begin{taxonomy} \hlt{Core Order-Type Segmentation}
\begin{enumerate}[label=\roman*.]
\setlength{\itemsep}{0pt}
\item \hlt{market}: pay immediacy and consume the best available liquidity;
\item \hlt{limit}: specify a worst acceptable price and wait in the book if not marketable;
\item \hlt{stop}: become active only after a trigger price is reached;
\item \hlt{IOC/FOK}: cancel any unfilled quantity immediately, or require full immediate execution;
\item \hlt{pegged}: reference the midpoint, best bid, best offer, or another venue-derived price;
\item \hlt{hidden/iceberg}: reduce displayed size while keeping some executable interest undisplayed.
\end{enumerate}
\end{taxonomy}

\begin{principle} \hlt{Order Choice Trades Certainty for Leakage}\\
Aggressive orders buy execution certainty but reveal urgency. Passive and hidden orders reduce signalling, but they introduce queue risk, missed fills, and adverse selection.
\end{principle}

\begin{caveat} \hlt{Instruction Semantics Are Venue-Specific}\\
The same label can behave differently across venues. Minimum quantity, display refresh, priority loss after modification, pegging reference, and cancellation rules must be read as part of the market design.
\end{caveat}
